\section{Performance Evaluation and Metrics}

To validate the viability of the DBaaS platform as a production-grade system, a series of performance evaluations were conducted to measure infrastructure efficiency, provisioning latency, and proxy overhead.

\subsection{Database Provisioning Latency}
Traditional managed database services (such as AWS RDS or Google Cloud SQL) typically require 5 to 10 minutes to allocate compute instances, configure networking, and bootstrap a replication cluster \cite{aws_rds_provisioning}. By leveraging the CloudNativePG operator and Kubernetes-native abstractions, the platform achieves extremely low latency in provisioning:
\begin{itemize}
    \item \textbf{Instantiation Speed:} From the moment a user initiates a "Create Project" request via the Flutter frontend, the Go backend generates the CRD, and the K8s operator schedules the underlying pods in an average of \textbf{35 to 45 seconds}.
    \item \textbf{Zero-Downtime Readiness:} This timeframe is fully inclusive; it encapsulates the generation of cryptographic secrets, the attachment of Persistent Volume Claims (PVCs), and the instantiation of the PostgREST auto-API sidecar. 
\end{itemize}

\subsection{PGProxy Routing Overhead}
Because the PGProxy intercepts all external database traffic, it was imperative to ensure it did not introduce unacceptable latency bottlenecks. To empirically validate the proxy's performance, the official PostgreSQL benchmarking tool (\texttt{pgbench}) was executed under two conditions: a direct intra-cluster connection to the database pod, and an external connection routed through the PGProxy.

\begin{itemize}
    \item \textbf{Benchmarking Results:} The \texttt{pgbench} evaluation established a direct intra-cluster baseline latency of \textbf{12.7 milliseconds} per transaction. When routed externally through the PGProxy, the system imposed an almost negligible latency penalty (measured between \textbf{1.5 and 3.0 milliseconds} of overhead) per transaction when compared to the direct connection baseline.
    \item \textbf{Zero-Copy Peeking:} This minimal overhead is achieved by implementing zero-copy TLS peeking on the \texttt{ClientHello} packet rather than full cryptographic termination, leveraging Go's highly optimized \texttt{net} and \texttt{io} standard libraries \cite{go_net_performance}.
    \item \textbf{Throughput Efficiency:} The proxy's streaming architecture ensures that once the initial SNI resolution is complete, bytes are transparently piped between the client and the K8s pod at wire speed, ensuring high throughput for large analytical queries or \texttt{pg\_dump} backup streams without degrading queries per second (QPS).
\end{itemize}

\subsection{Resource Footprint of Idle Tenants}
Multi-tenant density is a critical metric for DBaaS cost efficiency. To maximize the number of tenants per node, the base configuration of tenant databases is heavily optimized:
\begin{itemize}
    \item \textbf{Minimal Memory Allocation:} A standard "Free Tier" idle PostgreSQL tenant running via CloudNativePG, alongside its PostgREST sidecar, consumes an astonishingly low \textbf{65MB to 80MB} of RAM combined. This extreme efficiency allows a single 16GB RAM worker node to comfortably host over \textbf{150 idle database clusters} simultaneously.
    \item \textbf{Compute Quotas:} Strict Kubernetes CPU limits (e.g., \texttt{0.5 vCPU}) and memory quotas are enforced at the Namespace level, guaranteeing that a single noisy neighbor cannot monopolize the host cluster's resources and degrade performance for other users.
\end{itemize}
